\chapter{Discussion}
\label{ch:discussion}

\Cref{ch:results} reported what happened. This chapter asks what it means: what the numbers say
about the research questions, why the central mechanism works, which measurement decisions changed
the conclusions, and where the work's limits are.

\section{Answering the research questions}
\label{sec:answers}

\paragraph{RQ1 --- Does the neuro-symbolic model outperform deep reinforcement learning of equal
information?}
On score alone, the honest answer is \emph{barely}: \num{53.81} against \num{52.35} over all
scenarios, and on the six training scenarios the difference is not statistically supported
($p=0.213$). On the terms the problem actually sets, the answer is \emph{yes, and not narrowly}:
the \NS never violated a constraint in \num{360} shifts, while \texttt{ppo} did so in \num{53}, and
it reached its throughput with \SI{21}{\percent} less energy (\SI{8.42}{\watt\hour} against
\SI{10.21}{\watt\hour}) and a mean minimum battery of \SI{59.0}{\percent} against
\SI{45.0}{\percent}. Both policies solve a problem well; only one of them solves the problem as
stated in \cref{sec:constraints}.

\paragraph{RQ2 --- What does each half contribute?}
The rules contribute the guarantee, and they are not merely a safety net: the symbolic layer alone
scores \num{52.91}, above \texttt{ppo} and far above the rule-based reference. The network
contributes a small, consistent improvement on top --- \num{0.90} points against the symbolic layer
and \num{0.96} against the analytic estimate alone, both at $p<0.001$. Neither half is redundant,
but they are not equal partners: the rules carry most of the competence, and the network refines
it. The most striking evidence for the division of labour is \texttt{ns\_neural\_only}. Without
rules, the same network delivers \emph{more} units than the full model (\num{60.7} against
\num{59.9}) and loses almost no production (\num{0.21} against \num{1.90}) --- it has genuinely
learned what is worth doing --- while violating constraints in \num{57} of \num{360} shifts. It
learned the judgement and not the guarantee, which is precisely the split the architecture assumes.

\paragraph{RQ3 --- Do the policies hold up in situations they were never trained for?}
This is where the approaches separate most clearly. On the six unseen scenarios the \NS scores
\num{47.20}, its unbounded variant \num{43.33}, the symbolic layer \num{45.66} and \texttt{ppo}
\num{44.83} with \num{49} unsafe shifts out of \num{180}. Two observations follow. First, an
unbounded learned ranker is worse than no ranker at all outside its training distribution: the
network does not merely stop helping, it actively harms. Second, PPO's failures are concentrated
where its training no longer covers the situation --- \num{49} of its \num{53} unsafe shifts are in
unseen scenarios, \num{29} of them in \texttt{heavy\_load} alone.

\paragraph{RQ4 --- Do fast-simulator results transfer to the full simulation?}
Yes, with a known bias. Energy transfers almost exactly (\SI{+0.2}{\percent}) because it is
computed from the same physical model in both cases; throughput is \SI{5.1}{\percent} lower in
Gazebo because real driving loses seconds per leg to acceleration, arrival tolerances and
occasional recovery behaviours. The bias is systematic and in one direction, so it shifts all
policies together and leaves their ordering intact --- which is what a comparison needs. The
sixty-minute shifts additionally reproduced the safety difference itself
(\cref{tab:longshift}), so the central claim of this thesis does not rest on the abstraction.

\section{Why bounding the correction works}
\label{sec:why-bound}

The bound of \cref{eq:bounded} is one line of code, and it is worth being precise about why such a
small change is worth \num{3.87} points on unseen scenarios.

The score of an action is the sum of two terms with different epistemic status. The analytic
estimate is arithmetic on the current state: given the distances, the payload and the objective
weights, it computes what an action moves and what it costs. Arithmetic \emph{extrapolates} --- it
is as valid in a factory producing twice as fast as in the one it was written for. The learned
correction is a function fitted to a sample of states; it \emph{interpolates}. Inside the region
its training covered, it encodes real knowledge that the arithmetic lacks: which line will block
next, when a fuller load repays the detour. Outside that region, it continues to produce confident
numbers with no basis.

Unbounded, the second term dominates the first: the corrections observed in the failing scenario
were around $+20$ on every candidate, while the estimates differed by a few units. The ranking was
therefore decided entirely by the network, in a factory where its training did not apply, and the
rules could not intervene because none of its choices were illegal --- waiting with cargo while
lines block is wasteful, not forbidden.

Bounding the correction re-establishes the intended hierarchy without touching either component.
Where the arithmetic is nearly indifferent, the network decides, and its judgement is exactly what
is wanted. Where the arithmetic is clear, no amount of misplaced confidence can overturn it. The
selection table (\cref{tab:bound-sweep}) shows both halves of this: $\beta=0$, which silences the
network, is worse than every bounded setting, and so is $\beta=\infty$.

This is a different mechanism from the shielding and action masking discussed in
\cref{sec:safe-learning}. Those act on the \emph{legality} of actions: a shield removes what is
unsafe, and the learned policy is free within what remains. The bound acts on the \emph{influence}
of the learned component over choices that are all legal. The two are complementary --- this system
uses both --- and the distinction is worth stating, because a shield would not have prevented the
failure described above.

\section{Why the rules matter more than the score suggests}
\label{sec:rules-matter}

A reader comparing only scores would conclude that the difference between the approaches is minor.
The comparison of \texttt{ppo\_wp6} with the \NS makes the point sharply: the two are
statistically indistinguishable in score ($p=0.121$), and one of them enters an unsafe state in
\num{51} of \num{360} shifts.

What a violation means physically is not a slightly lower score. The battery reserve exists so the
robot can always reach the charging station; below it, the robot risks stopping in a corridor,
where it blocks a route, requires a person to fetch it, and delivers nothing further that shift.
The scoring penalty of \num{20} points is a modelling convenience, not a real cost --- the real
cost is an unplanned human intervention in a factory that was supposed to run by itself. This is
why \cref{ch:results} reports unsafe shifts beside every score rather than trusting the penalty to
express the difference.

There is a second, quieter argument for the rules. The \NS explains every decision in terms a
plant engineer can check: which constraint forbade an alternative, which priority applied, by how
much the network preferred the chosen action. When the robot behaved oddly in
\cref{sec:hard-rules}, that explanation is what made the cause findable in minutes. A policy whose
output is an action index affords no such inspection --- a property that matters more in an
industrial setting than a fraction of a point of score.

\section{Measurement decisions that changed the conclusions}
\label{sec:measurement}

Two results in this thesis exist only because a measurement was questioned, and both changed a
decision that had already been made.

\paragraph{An unfair training setup.}
PPO's parallel environments were originally each bound to one scenario, so with four environments
the agent saw four of the six training scenarios. Had this gone unnoticed, the thesis would have
compared a fully trained neuro-symbolic model against a partially trained baseline, and the
comparison would have been worthless regardless of its statistics. The correction is undramatic
--- after retraining on all six, PPO's score and its violation rate are essentially unchanged ---
but that outcome is itself the useful finding: PPO's constraint violations are not a symptom of
insufficient scenario coverage, and would not have been fixed by more of the same training.

\paragraph{A benchmark that measured the wrong regime.}
The global planner was selected on a short benchmark tour, where Theta$^{*}$ was \SI{3.3}{\percent}
faster with half the turning --- a clean, statistically significant result ($p=0.001$). Over full
shifts the same planner produced faults in \num{6} of \num{17} shifts, including one that ended
with the robot halted, because it refuses to plan when the start or goal cell is momentarily
occupied. The benchmark was not wrong; it answered the question it was asked, which was about
tours between free poses in an empty hall. The operating regime --- a robot standing against a
machine, in a costmap containing fresh sensor data --- was outside it. The lesson generalises
beyond planners: a benchmark measures a regime, and a component chosen in one regime has to be
re-checked in the one it will run in.

\section{Limitations}
\label{sec:limitations}

\paragraph{The objective does not charge for charging.}
Only energy drawn from the battery is counted (\cref{sec:objective}). A robot parked on the charger
draws its idle power from the charging station, so waiting there is free in the score although the
factory pays for it. This favours policies that spend time docked, and it is the main reason the
rule-based reference appears so energy-efficient (\SI{0.099}{\watt\hour} per unit against
\SI{0.150}{\watt\hour}): it charges \num{7.1} times per shift and spends much of the hour parked.
Correcting this would require re-training both learned models against the corrected objective, so
it is reported rather than fixed. The comparison between policies remains valid --- all of them are
scored the same way --- but the absolute energy figures understate what a real installation would
pay.

\paragraph{The remaining weakness of the model is a batching trade-off.}
In \texttt{worn\_robot}, the only scenario where a simpler policy wins significantly, the rule
reference makes \num{14} deliveries of \num{3.7} units and drives \SI{359}{\metre}; the \NS makes
\num{25} deliveries of \num{2.2} units and drives \SI{468}{\metre} for the same output. Reacting
quickly is what wins everywhere else, and it is what loses here, because a worn drivetrain makes
every metre more expensive. The model has no notion that waiting for a fuller load can be the
cheaper option; its look-ahead horizon of \SI{900}{\second} may simply be too short for that
trade-off to appear in the training targets.

\paragraph{One robot, one hall, one objective.}
The problem has a single robot, a single delivery point and a single charger, and no human traffic.
Several robots would introduce congestion, charger contention and task allocation --- a different
problem class in which the symbolic layer's per-action reasoning would have to be extended
considerably. Likewise, all results come from one maze layout and one set of objective weights;
the ratio between the cost of a lost unit and the cost of a watt-hour in particular is a modelling
choice, and a factory that valued energy differently might rank the policies differently.

\paragraph{The symbolic layer is hand-written.}
Its rules are only as good as the understanding of whoever wrote them, and H5 is direct evidence:
a necessary rule was missing until a real shift exposed it. There is no mechanism guaranteeing that
no further such gap exists. Learning the symbolic layer, as some of the related work does
(\cref{sec:decision-making}), would address this, at the cost of the readability the hand-written
rules provide.

\paragraph{Simulation throughout.}
No physical robot was involved. The energy model is calibrated and validated against Gazebo, but
Gazebo shares its assumptions about rolling resistance and motor efficiency, so agreement between
them is consistency rather than proof. Localisation error, wheel slip, floor variation and battery
ageing would all test the model in ways this work cannot.

\paragraph{The full-simulation evidence is thin in volume.}
Twelve twenty-minute shifts and two sixty-minute shifts are demonstrations, not statistics. They
establish that the fast simulator does not mislead and that the safety difference is real in the
full stack; they cannot support a confidence interval on the sim-to-real gap.

\section{Threats to validity}
\label{sec:validity}

\paragraph{Internal validity.}
Policies are compared on identical seeds with paired non-parametric tests, which removes the
dominant source of variance. No correction for multiple comparisons is applied
(\cref{sec:statistics}), so isolated significant results should be treated with caution; the
conclusions drawn here rest on differences that are consistent across scenarios or supported by a
mechanism, not on single $p$-values. The bound $\beta$ was selected on training scenarios and
development seeds only, so the improvement it produces on the test set is not a selection effect.

\paragraph{External validity.}
The results describe this factory: one layout, three sections with these rates, this robot, this
objective. The scenario set was deliberately widened to twelve configurations, six of them never
seen during training, which supports claims about robustness to \emph{parameter} changes --- rates,
masses, buffer sizes, battery health, ambient temperature. It supports no claim about a different
hall geometry, more robots, or a different task structure.

\paragraph{Construct validity.}
Whether the score in \cref{eq:objective} expresses what a factory wants is a modelling assumption.
Two of its weights are defensible from first principles (a delivered unit and a lost unit are
worth the same; energy is cheap relative to production), while the penalty of \num{20} per safety
violation is a stand-in for a cost --- an unplanned intervention --- that has no natural unit. This
thesis therefore avoids resting any conclusion on the penalty's magnitude, and reports violation
counts separately throughout.
